\documentclass[12pt]{article}

\usepackage[utf8]{inputenc}
\usepackage[T1]{fontenc}

\title{La linguistica de la biologia}
\author{Dana Barandica,Junior Sabogal, Dulce Velasquez y Miguel Lopez}
\date{\today}

\begin{document}

\maketitle

\section{Introducción}
La linguistica de la biologia analiza el lenguaje como una propiedad biologica propia de nosotros los humanos# realmente no es propia de los humanos# y no unicamente algo de costumbre social o cultural# NO ES CLARO. Su objeto principal de estudio # CUÁL ES EL SUJETO es la Lengua-I (interna y/o individual)#?, que es un sistema de conocimiento sobre el lenguaje que desarrolla la mente y cerebro de cada persona #ANTI-REALISMO?. Tras la teoria de la evolucion de Darwin, surgieron los primeros analisis que comparaban la evolucion gradual de las especies con la de las lenguas.#EN ESTE PARRAFO HAY VARIAS IDEAS QUE NO ESTAN CLARAMENTE ARTICULADAS.

Los seres humanos somos las unicas especies conocidas #ESTO NO ES CORRECTO que presenta la capacidad de comunicarse de diferentes maneras, tantas formas hay de comunicarnos #REDUNDANCIA y aun asi somos seres tan organicos#? que no somos capaces de entender el sufrimiento de los demas# NO ES CLARA ESTA IDEA. Aun asi a pesar que tenemos tantas formas de comunicarnos la ciencia encontro otra manera, mediante la obtencion de datos por medio de un mecanismo, esta manera es un reglamento universal que todos los seres humanos que quieren escribir ciencia deben seguir por lo menos #LO ANTERIOR NO ES CLARO.


Es diferente el lenjuaje humano y los diversos sistemas de señales biologicas. La biocomunicacion # ATRÁS MENCIONAN QUE SOLO LOS SERES HUMANSO SE COMUNICAN, NO ES CLARO abarca todas las interraciones ya sea intra e interespecificas en plantas, animales, hongos y microorganismos mediante señales quimicas basadas en diferentes moleculas, y las vocalizaciones que siguen reglas comportamentales. el lenjuaje humano no evoluciono principalmente por presiones selectivas de comunicacion, sino para simbolizar el moldeamiento del pensamiento.#IDEM

IDEAS PRINCIPALES:
BIOLINGÜISTICA: Se analizará el lenguaje como una capacidad biológica, de forma que no se infiera la lingüistica solo como un ámbito sociocultural, si no como una facultad natural de la humanidad, permitiendo la existencia de una " Gramática Universal", la cual se podria comparar con un código base o un modelo mental compartido, donde no solo se atribuye a un rasgo genético o estrechamente biológico, si no a códigos lingüisticos como (la expresion de pensamientos, el simbolismo, la interpretación y la sintaxis), viendo a la lingüistica no como añadidura, si no como un complemento natural que construye la comunicación y a un orden grabado de la esencia innata del ser humano.
BIOSEMIÓTICA: Argumentar que la funcion primaria del lenguaje pudo ser la respresentacion simbolica y el pensamiento productivo abstracto, plateando a la comunicacion como una funcion secundaria. Fenómeno simbólico donde la interpretacion de signos y codigos es inmanente a todo ser vivo, desde la unidad más pequeña del ser humano (celula), hasta el humano.
PARALELISMO ENTRE CODIGO GENÉTICO Y LINGÜISTICO: Abordar que el código genetico prefigura el código lingüistico, planteando que ambos son sistemas de transmisión de información con aspectos parecidos. Al igual que el lenguaje humano puede formar multiples palabras y mensajes con pocos recursos, el código genético puede producir millones de anticuerpos para neutralizar antÍgenos desconocidos.
i
\end{document}
